% 15_body.tex — Day3 산출물 5/5 (⑮) 조달·계약 변경 기록 — 변경 대가 · 수행사 성과 · 검수 보류 회신 (빈 템플릿 / 완성 예시 공용 본문)
% 드라이버: 15_조달계약변경_빈템플릿.tex (\wbblanktrue) / 15_조달계약변경_완성예시.tex (\wbblankfalse)
% 내용 출처: PM트랙/Day3_워크북.md §5.3(항목 가이드) §5.4(빈 템플릿) §5.5(온담 P1 완성 예시본)
\documentclass[10pt,a4paper]{article}
\usepackage[a4paper,margin=16mm,top=14mm,bottom=20mm]{geometry}
\def\DocNum{5}
\def\DocTitle{조달 · 계약 변경 기록}
\def\DocSub{⑮ Procurement \& Contract Change Record}
\def\DocVariant{\B{빈 템플릿}{온담 P1 완성 예시}}
\usepackage{wbstyle}
\begin{document}

\docheader
 {\Bg{[정식 명칭과 코드]}{차세대 주문·재고 통합 플랫폼 구축 (ONE ONDAM P1)}}
 {\Bg{[기간 / 데이터 일자 / 발행]}{2026-01-05 ~ 09-30 / 발행 2026-10-05}}
 {\Bg{[P1 PM · 구매]}{P1 PM · 강현도 구매팀장\rd{5mm}}}
 {\Bg{[스폰서 / 재경]}{정미란 (5억 초과 병행)\rd{5mm}}}

% ------------------------------------------------------------------ 1
\secbar{1. 문서 정보와 계약 관리 역할}
\guide{기록 기간·데이터 일자·발행, 부속 관계(성과보고·변경협의체 자료), 계약 관리 역할 — 누가 변경협의체의 발주 대표이고, 누가 서면·단가·조달을, 누가 지급·내부통제를, 누가 하도급·위탁 검토를 맡는가. 검수 조항의 시한(15영업일)과 변경 서면의 번호 체계를 PM이 알고 있어야 한다.}
{\footnotesize
\begin{tabularx}{\linewidth}{|L{40mm}|Y|}\hline
\hd{항목} & \hd{내용} \\ \hline
\B{\rd{9mm}}{}기간 / 데이터 일자 / 발행 & \Bg{[계약 착수 ~ 데이터 일자 / 발행, 부속 관계]}{2026-01-05 ~ 09-30 / 09-30 / 10-05 (성과보고 제9호 부속, 변경협의체 10-14 자료)} \\ \hline
\B{\rd{16mm}}{}계약 관리 역할 & \Bg{[PM / 구매 / 재경 / 법무 / PMO — 각자의 결정 권한과 서명 위치]}{P1 PM(변경협의체 발주 대표, 검수 결과 통보), 강현도 구매팀장(계약 서면·단가·조달), 재경팀장(지급·내부통제 5,000만 원 이상 CFO 결재), 법무(최영주 — 하도급·위탁 검토), P1 PMO 기준선 담당(FP 확인 3인 중 1)} \\ \hline
\B{\rd{12mm}}{}관리 대상 & \Bg{[P1 예산 내 계약 n건 + P1 산출물 의존 계약 n건 — 소관 명시]}{P1 예산 내 9건(SI, API GW, MDM, APM/SAP 애드온, 클라우드, 데이터 전환 용역, 감리, 보안진단, POS 단말 예정) + P1 산출물 의존 2건(자사앱 벤더 — 영업본부, 3PL — 물류본부). 상세는 1-1.} \\ \hline
\B{\rd{10mm}}{}변경 서면 번호 체계 & \Bg{[계약별 접두어 — 로그(항목 7)와 같은 번호]}{SI 변경 합의서 n호 / SW-nn(상용SW) / CL-nn(클라우드) / DT-nn(전환 용역) / AU-nn(감리) / 하도급 승인서 n호 / 앱-nn·3PL 부속(P1 예산 외)} \\ \hline
\B{\rd{8mm}}{}작성 / 검토 / 확인 & \Bg{[ ]}{P1 PM / 강현도(계약 서면 대조), 재경팀장(금액) / 정미란(5억 초과 건)} \\ \hline
\end{tabularx}}

% ------------------------------------------------------------------ 1-1 (가로)
\landscapeon
\secbar{1-1. 계약 구조 — 관리 대상 계약 목록}
\guide{계약명·상대·금액·형태(고정가/단가/종량)·소관과 P1 예산 여부·검수 조건(조 번호)·지체상금(조 번호·율·상한)·변경 절차·서면 번호. P1 예산 밖 계약도 “P1 산출물 의존”으로 포함하되 소관을 적는다 — SI 계약만 관리하면 앱 벤더·3PL 지연이 어느 문서에도 없다.}
{\scriptsize\setlength{\tabcolsep}{2.5pt}
\begin{tabularx}{\linewidth}{|L{26mm}|L{26mm}|L{38mm}|L{30mm}|Y|L{34mm}|L{36mm}|L{18mm}|}\hline
\hd{계약} & \hd{상대} & \hd{금액 (억) · 형태} & \hd{소관 · P1 예산} & \hd{검수 조항} & \hd{지체상금} & \hd{변경 절차} & \hd{서면 번호} \\ \hline
\ifwbblank
\rd{11mm}\g{[SI]} & & \g{[고정가 FP n × 단가]} & & \g{[§n 릴리스별, n영업일, 발주사 사유 조항]} & \g{[§n 율/일, 상한]} & \g{[§n 협의체 → 합의서]} & \\ \hline
\rd{9mm}\g{[상용SW]} & & & & \g{[설치 확인서 / 통합시험 통과]} & & & \\ \hline
\rd{9mm} & & & & & & & \\ \hline
\rd{9mm}\g{[클라우드]} & & \g{[선약정 + 종량]} & & & & & \\ \hline
\rd{9mm}\g{[용역]} & & & & & & & \\ \hline
\rd{9mm}\g{[감리 / 진단]} & & & & & & & \\ \hline
\rd{9mm}\g{[물품 (예정)]} & & \g{[기준선 → 견적]} & & \g{[입고 검수]} & & \g{[발주 일자]} & \\ \hline
\rd{9mm}\g{[P1 예산 외 의존 1]} & & & \g{[소관 본부 — P1 의존 WBS]} & & & & \\ \hline
\rd{9mm}\g{[P1 예산 외 의존 2]} & & & & & & & \\ \hline
\else
SI 개발 용역 & ㈜한빛디지털 (이재현) & 68.20, 고정가 FP 12,400 × 55만 원 & P1 (IT본부) & §14 릴리스별 검수 요청 후 15영업일 내 결과 통보, 발주사 사유 지연 기간은 지체 제외·대기 비용 대상 & §18 지연 부분 계약금 × 1/1000/일, 상한 10\% & §12 변경협의체(격주 수) 합의 → CCB → 변경 합의서 & SI 합의서 n호 \\ \hline
API 게이트웨이 & SW 벤더 A & 라이선스 5.6 + 유지보수, 인도 기준 & P1 & 설치 확인서 & — & SW-nn & SW-01, 04 \\ \hline
상품·고객 MDM & SW 벤더 B & 7.6 + 유지보수 & P1 & 설치 확인서 & — & SW-nn & SW-02, 04 \\ \hline
APM·모니터링 / SAP 애드온 & SW 벤더 C / SAP 파트너 & 3.4 / 2.8 + 유지보수 & P1 & 설치 확인서 & — & SW-nn & SW-03, 04 \\ \hline
클라우드 & CSP & 2년 선약정 6.0 + 종량 (기준선 11.6) & P1 & 월 사용량 청구 & — & CL-nn & CL-01 \\ \hline
데이터 전환 용역 & 전환 용역사 & 11.4 → 9.4 (고정가, 월 청구) & P1 & 리허설·본 전환 검증 6종 & §9 1/1000/일 & DT-nn & DT-01 (예정) \\ \hline
외부 감리 & 감리법인 (정하늘) & 4.2 (G1·중간 2회·G3·G4) & P1 & 회차 보고서 & — & AU-nn & AU-01 \\ \hline
보안취약점 진단 & 3자 진단 & 2.4 (2회) & P1 & 진단 보고서 & — & — & — \\ \hline
POS 단말 (예정) & 단말 공급사 & 8.6 → 견적 8.9 & P1 (강현도) & 입고 검수·설치 확인 & 물품 1/1000 & 발주 10-30 & — \\ \hline
\emph{자사앱 유지보수} & 앱 외주 벤더 & 월 3,200만 원 & \textbf{영업본부 (P1 예산 외)}, P1 산출물 의존(1.2.3) & — & — & 영업본부 & 앱-01 \\ \hline
\emph{3PL} & ㈜대성로지스 & 3PL 수수료 & \textbf{물류본부 (P1 예산 외)}, P1 의존(1.3.3) & — & — & 물류본부 & 3PL 부속 \\ \hline
\fi
\end{tabularx}}

% ------------------------------------------------------------------ 2
\secbar[70mm]{2. 변경 대가 산정표 \B{}{— SI 계약, 단가 550,000원/FP}}
\guide{CR별 FP 델타 × 계약 단가 = 금액(증액은 +, 감액은 −), 대기 비용(판정 일수 × 대기 인력 × 일당 — 산식은 계약 조항으로), 변경협의체 합의 일자, 계약 변경 서면 번호·서명 일자, 지급 시점(검수 연동), 누적. FP 산정은 수행사 산정서 + 발주사 PMO 검증(FP 확인 3인). 변경 대가를 종료 시 일괄 협상하지 않는다 — 그때는 근거 문서가 없다. 감액을 계약에 반영하지 않으면 범위는 줄었는데 대금은 그대로다.}
{\footnotesize\setlength{\tabcolsep}{2.5pt}
\begin{tabularx}{\linewidth}{|L{40mm}|C{13mm}|C{16mm}|C{12mm}|Y|L{34mm}|L{36mm}|L{34mm}|}\hline
\hd{CR / 건} & \hd{FP Δ} & \hd{금액 (억)} & \hd{증/감} & \hd{대기 비용 (일수 × 인력 × 일당, 산식 조항)} & \hd{변경협의체 합의} & \hd{서면 (서명일)} & \hd{지급 시점} \\ \hline
\ifwbblank
\rd{8mm}CR- & & \g{[Δ × 단가]} & & & & & \g{[릴리스 검수 후]} \\ \hline
\rd{8mm}CR- & & & & & & & \\ \hline
\rd{8mm}\g{[대기 비용 (IS-)]} & — & & 대기 & \g{[n일 × n명 × 일당 — §n]} & \g{[판정: n일 중 발주사 귀책 n일]} & & \\ \hline
\rd{8mm}CR- & & & & & & & \\ \hline
\rd{8mm}CR- & & & & & & & \\ \hline
\rd{8mm}CR- & & & & & & & \\ \hline
\rd{8mm}CR- & & & \g{[감액은 −]} & & & & \g{[잔금 차감]} \\ \hline
\rd{8mm}\textbf{누적 (SI)} & \g{[±n]} & \g{[ ]} & & \g{[ ]} & & & \g{[합]} \\ \hline
\rd{8mm}\g{[타 계약 감액·증액]} & — & & & & & & \g{[기준선 반영 여부]} \\ \hline
\else
CR-03 엑셀 양식 5종 & +22 & +0.121 & 증액 & — & 03-18 & SI 합의서 3호 (03-25) & R2 검수 후 (10-16) \\ \hline
CR-06 FZ-05 승인 16건 & +58 & +0.319 & 증액 & — & 06-24 & SI 합의서 5호 (07-15) & R3 검수 후 \\ \hline
R-05 대기 비용 (IS-02) & — & +0.42 & 대기 & 7영업일 × 6명 × 100만 원/일 [추가 설정 — §14 산식: 대기 인력 × 계약 단가 환산 일당] & 07-15 (판정: 12일 중 발주사 귀책 7일) & SI 합의서 5호 & 07-31 지급 완료 \\ \hline
CR-07 인터페이스 2본 & +40 & +0.22 & 증액 & — & 09-16 & SI 합의서 6호 (10-23 예정) & R3 검수 후 \\ \hline
CR-08 파일럿 임시 연동 & +90 & +0.495 & 증액 & — & 09-30 & 6호 & R4 검수 후 \\ \hline
CR-09 22:00 파라미터 & +35 & +0.1925 & 증액 & — & 09-30 & 6호 & R3 검수 후 \\ \hline
CR-11 5년 초과 이력 제외 & −60 & −0.33 & \textbf{감액} & — & 09-30 (수행사 이의 없음) & 6호 & 잔금에서 차감 \\ \hline
\textbf{누적 (SI)} & \textbf{+185} & \textbf{+1.0175} & 순증 & \textbf{+0.42} & & & \textbf{합 1.4375} \\ \hline
데이터 전환 용역 (CR-11) & — & \textbf{−2.00} & 감액 & — & 09-25 (용역사 합의, 강현도) & DT-01 (10-23 예정) & 11.4 → 9.4 (BL-01 기반영) \\ \hline
\fi
\end{tabularx}}
\B{\small\g{FP 확인 3인 서명 현황과 대기 비용 일당 산식(조항·단가표): \hrulefill}\par}{\small FP 확인 3인(P1 PMO 기준선 담당·이재현·형상 담당) 서명: CR-03·06 완료, CR-07~11은 BL-02 FP 확인서(10-23)에서 일괄. 대기 비용 일당 산식은 §14의 “대기 인력의 계약 단가 환산 일당”을 P-11a(44명 × 640만 원/월 ÷ 20영업일 = 인당 약 32만 원/일)에 간접비·이윤을 얹은 계약 부속서 2의 단가표(인당 100만 원/일 [추가 설정])로 산정. 검산: 22+58+40+90+35−60 = 185; 185 × 0.0055 = 1.0175; 1.0175 + 0.42 = 1.4375.\par}

% ------------------------------------------------------------------ 3
\clearpage
\secbar{3. 수행사 성과 기록 \B{}{— 스프린트 S7~S18 (구현 12회)}}
\guide{스프린트별 계획·인수 FP(velocity)와 \textbf{이월 누적}(계획 대비 \%만 적으면 이월이 안 보인다), 결함, 인력 계획/실적, 핵심 인력 교체, \textbf{발주사 결정 SLA 준수(양방향)}. 성과는 사실만 — 평가는 종료 보고서(Day4)에서. 발주사 측 의무 준수를 함께 적지 않은 일방적 기록은 협상에서 쓸 수 없다.}
{\footnotesize\setlength{\tabcolsep}{2.5pt}
\begin{tabularx}{\linewidth}{|L{42mm}|C{16mm}|C{18mm}|C{16mm}|L{34mm}|L{44mm}|Y|}\hline
\hd{스프린트 (실제 기간)} & \hd{계획 FP} & \hd{인수 FP} & \hd{이월 누적} & \hd{결함 (스프린트 기간)} & \hd{인력 계획/실 (명)} & \hd{발주사 결정 SLA 준수 (5영업일 내 응답률)} \\ \hline
\ifwbblank
\rd{5mm}S & & & & & & \\ \hline
\rd{5mm}S & & & & & & \\ \hline
\rd{5mm}S & & & & & & \\ \hline
\rd{5mm}S & & & & & & \\ \hline
\rd{5mm}S & & & & & & \\ \hline
\rd{5mm}S & & & & & & \\ \hline
\rd{5mm}S & & & & & & \\ \hline
\rd{5mm}S & & & & & & \\ \hline
\rd{5mm}S & & & & & & \\ \hline
\rd{5mm}S & & & & & & \\ \hline
\rd{5mm}S & & & & & & \\ \hline
\rd{5mm}S & & & & & & \\ \hline
\rd{5mm}\textbf{합계 · 평균} & & \g{[n (n\%)]} & & \g{[n (n/FP)]} & \g{[평균 / 피크 계획]} & \g{[최저 월 → 최근 월]} \\ \hline
\else
S7 (03-30~04-12) & 620 & 620 & 0 & 36 & 46/46 & 100\% \\ \hline
S8 (04-13~04-26) & 620 & 620 & 0 & 69 & 48/48 & 92\% \\ \hline
S9 (04-27~05-10) & 620 & 600 & 20 & 77 & 50/49 & 83\% \\ \hline
\textbf{S10 (05-11~05-29, 3주)} & 620 & \textbf{480} & 160 & 120 & 50/50 & \textbf{58\%} (POS 화면 12영업일) \\ \hline
S11 (06-01~06-14) & 620 & 470 & 310 & 98 & 52/51 & 75\% \\ \hline
S12 (06-15~06-28) & 620 & 560 & 370 & 98 & 52/52 & 88\% \\ \hline
S13 (06-29~07-12) & 620 & 560 & 430 & 83 & 54/53 & 90\% \\ \hline
S14 (07-13~07-26) & 620 & 600 & 450 & 90 & 54/52 (PL 교체 인수인계) & 92\% \\ \hline
S15 (07-27~08-09) & 620 & 620 & 450 & 50 & 54/54 & 95\% \\ \hline
S16 (08-10~08-23) & 620 & 620 & 450 & 20 (시험 인력 규제 이동) & 54/54 & 94\% \\ \hline
S17 (08-24~09-06) & 620 & 640 & 430 & 71 & 54/56 (초과 근무) & 91\% \\ \hline
\textbf{S18 (09-07~09-25, 3주)} & 620 & 630 & \textbf{420} & 118 & 54/52 & 91\% \\ \hline
\textbf{합계 · 평균} & 7,440 & \textbf{7,020} (94.4\%) & 420 & 930 (0.132/FP) & 평균 52.3 / 피크 계획 61 (R3–R4) & 5월 58\% → 9월 91\% \\ \hline
\fi
\end{tabularx}}
\B{\small\g{계약상 인도물 준수(리뷰·FP 확인서·소스 동기화·설계문서 양식·시험 보고)와 결정 요청 SLA 총계: \hrulefill}\lines{2}{6mm}}{\small 계약상 인도물 준수: 스프린트 리뷰 12/12 실시, 인수 FP 확인서 12/12, 소스 동기화 주 1회 → 감리 7번 후 일 1회(10-02~), 설계문서 47본 중 38본 표준 양식 충족(9본 보완 10-30), 시험 결과 보고 12/12. 결정 요청 SLA: 수행사 → 발주사 요청 총 87건 중 5영업일 내 응답 76건(87\%), 초과 11건(5·6월 집중). 검산: 인수 620+620+600+480+470+560+560+600+620+620+640+630 = 7,020; 7,440 − 7,020 = 420.\par}
\landscapeoff

% ------------------------------------------------------------------ 3-1 인력 교체 승인서 (세로)
\secbar[60mm]{3-1. 핵심 인력 교체 승인서}
\guide{통보 일자·대상(역할과 담당 WBS)·사유, 후임(경력·증빙 검토자), 승인 조건(동시 투입 기간과 비용 부담·지식 이전 체크리스트·후임 주관 시험·향후 통보 기간), 승인·서명, 사후 확인. 인력 교체를 구두로 승인하지 않는다 — 승인 조건이 없으면 인수인계 기간의 속도 저하가 발주사 귀책으로 남는다.}
{\footnotesize
\begin{tabularx}{\linewidth}{|L{30mm}|Y|}\hline
\hd{항목} & \hd{내용} \\ \hline
\B{\rd{10mm}}{}통보 / 대상 & \Bg{[일자, 역할(담당 WBS), 사유, 퇴직·이동 예정일]}{2026-07-20, 인터페이스 PL(1.3.3 3PL 연동·1.5.5 47본 재구축 리더) 개인 사유 퇴직(08-14 예정)} \\ \hline
\B{\rd{10mm}}{}후임 & \Bg{[경력·유사 실적, 증빙 검토 일자·검토자]}{한빛디지털 인터페이스 아키텍트(경력 14년, SAP·EDI 연동 프로젝트 3건) — 이력서·경력 증빙 07-22 검토(박세연·P1 PM)} \\ \hline
\B{\rd{18mm}}{}승인 조건 & \Bg{[① 동시 투입 기간·비용 부담 ② 지식 이전 체크리스트(항목 수·확인자) ③ 후임 주관 시험 ④ 향후 통보 기간]}{① 2주 인수인계(07-20~08-02) 기간 전임·후임 동시 투입(비용 수행사 부담) ② 3PL 연동 명세 v1.0·13본 복원 지식 이전 체크리스트 24항목 확인(박세연 서명 08-03) ③ 폴백 시험(08-21) 후임 주관으로 통과 확인 ④ 이후 핵심 인력(PL 6명) 교체는 4주 전 통보} \\ \hline
\B{\rd{8mm}}{}승인 / 서명 & \Bg{[P1 PM / 수행사 PM / 기술 확인자 — 일자]}{P1 PM 07-23, 이재현 07-23, 박세연(기술 확인) 08-03} \\ \hline
\B{\rd{8mm}}{}사후 확인 & \Bg{[시험 결과, 후임 담당 건]}{폴백 시험 08-21 통과, 감리 2번(문서 양식 9본) 후임 담당} \\ \hline
\end{tabularx}}

% ------------------------------------------------------------------ 4 검수 보류 서면 회신
\secbar[80mm]{4. 검수 요청에 대한 서면 회신 — 진행 상태 · 완료 조건 · 발주사 사유 기간}
\guide{수신·참조·발신·일자·제목(수행사 문서 번호 인용) / ① 검수 요청 접수와 기산일·통보 기한(조항) ② 구두 통보의 정정 — “보류”가 아니라 진행 중 ③ 검수 완료 조건(수행사 조치·기한·조건부 합격의 대금 효과) ④ \textbf{발주사 사유 기간의 명시}(무엇·며칠·영향 범위 한정) ⑤ 계약상 효과(지체 제외·대기 비용 판정 절차) ⑥ 요청 사항 / 첨부. 발주사 사유를 0으로 적지 않는다 — 숨기면 변경협의체에서 수행사 기록이 이긴다.}
\begin{fieldbox}\small
\ifwbblank
\g{수신} \hrulefill\quad \g{참조} \hrulefill\par
\g{발신} \hrulefill\quad \g{일자} \hrulefill\par
\g{제목} \hrulefill\par
\g{1. 검수 요청의 접수 — 요청 일자·문서 번호, 대상 범위(WBS·FP), 기산일, 조항에 따른 결과 통보 기한}\lines{3}{7mm}
\g{2. 구두 통보의 정정 — 언제 누가 무엇을 말했고, 그것이 “보류”가 아님을 명시. 수행사 서면의 주장 인용}\lines{3}{7mm}
\g{3. 검수 완료 조건 — 수행사 조치 사항(감리 귀속 건·결함·RTM)과 기한, 조건부 합격의 대금 효과(조항)}\lines{4}{7mm}
\g{4. 발주사 사유 기간 — 무엇이 발주사 사유인지, 시작~종료(영업일), 영향받는 산출물의 범위(FP 비율), 병행 가능한 부분은 대기가 아님, 대기 기록 제출 기한과 판정 회의}\lines{5}{7mm}
\g{5. 계약상 효과 — 지체 일수 제외(조항), 대기 비용 반영 서면, 수행사 귀속 지연의 지체 산입}\lines{2}{7mm}
\g{6. 요청 사항 — ① 표현 정정 확인 ② 대기 기록 제출 ③ 조치 계획서 제출 (기한)}\lines{2}{7mm}
\g{첨부: }\lines{1}{7mm}
\else
\textbf{수신} ㈜한빛디지털 이재현 PM \quad \textbf{참조} 강현도 구매팀장, 재경팀장, 정하늘 수석감리원\par
\textbf{발신} ONE ONDAM P1 PM \quad \textbf{일자} 2026-10-01\par
\textbf{제목} 릴리스 R2 검수 요청(2026-09-25, 귀사 문서 HB-P1-2026-071)에 대한 회신 — 검수 진행 상태, 완료 조건 및 발주사 사유 기간의 명시\par
\textbf{1. 검수 요청의 접수.} 2026-09-25 귀사가 요청한 릴리스 R2 검수(대상: 재고 원장 코어 1.3.2, POS 거래·운영 코어 1.1.2·1.1.3, 결제·PG 인터페이스 7본 1.1.4, SAP 애드온·인터페이스 8본 1.5.4 일부, 규제 요건 반영분 — 인수 FP 누적 7,020 중 R2분 3,300 FP)을 접수하였으며, 계약서 §14에 따른 검수 기간 15영업일의 기산일은 2026-09-25, 결과 통보 기한은 \textbf{2026-10-16}입니다.\par
\textbf{2. 구두 통보의 정정.} 2026-09-28 주간 PM 회의에서 당사 PMO가 “감리 지적 정리 후 검수”라고 구두로 말한 것은 검수의 \textbf{보류}를 뜻하지 않습니다. 검수는 §14에 따라 진행 중이며, 귀사 09-29 서면(HB-P1-2026-073)의 “발주사 사유에 의한 검수 보류로 대기 비용 산정 대상”이라는 주장에는 아래 3·4항과 같이 회신합니다.\par
\textbf{3. 검수 완료 조건.} R2 검수 결과는 10-16에 통보하되, 다음 조건이 충족되지 않은 부분은 “조건부 합격(보완 요구)”으로 통보합니다. ① 중간감리 1차 지적 중 귀사 귀속 3건(AF-1-02 설계문서 9본 양식, AF-1-05 심각도 기준서 — 완료, AF-1-06 파기 주기 증적)의 조치 — 기한 10-16(06)·10-30(02) ② 심각도 2 미해결 결함 중 R2 관련 5건(D-0871·0902·0915·0944·0958)의 해소 — 기한 10-09~10-23 ③ RTM 검수 열의 R2 인수분 갱신. 조건부 합격 부분의 대금은 보완 확인 후 지급합니다(§15).\par
\textbf{4. 발주사 사유 기간의 명시.} 중간감리 1차 지적 AF-1-03(P2 인터페이스 규격 확인 서명)·AF-1-04(가맹 SW 배포 범위)는 당사가 2026-09-29 감리법인에 귀속 재분류를 요청한 건이며, 원인은 발주사(프로그램 PMO장·영업본부)의 결정 미완입니다. 따라서 \textbf{2026-09-28부터 귀속 정리 회신 예정일 10-08까지 8영업일은 발주사 사유 기간으로 인정}하며, 이 기간 중 해당 지적과 직접 관련된 산출물(P2 인터페이스 규격서 v1.1 갱신, 가맹 배포 패키지 설계 착수 — R2 FP의 약 12\%)에 대해 귀사가 작업을 진행하지 못한 사실은 §14의 “발주사 사유 지연 기간”에 해당함을 확인합니다. 다만 R2 검수 대상의 나머지 부분과 R3(S19) 작업은 병행 가능하므로 대기에 해당하지 않으며, 대기 인력·일수의 산정은 귀사가 10-09까지 제출하는 대기 기록(인력·일자·작업 불가 사유)을 근거로 \textbf{2026-10-14 변경협의체}에서 판정합니다.\par
\textbf{5. 계약상 효과.} 4항의 발주사 사유 기간은 R2 지체 일수 산정에서 제외되며(§14), 대기 비용은 변경협의체 판정에 따라 SI 변경 합의서 6호에 반영합니다. 3항의 귀사 귀속 조치 지연으로 인한 기간은 §18의 지체 일수에 포함됩니다.\par
\textbf{6. 요청 사항.} ① 09-29 서면의 “검수 보류” 표현을 본 회신에 따라 “검수 진행 중, 조건부”로 정정 확인 ② 대기 기록 10-09 제출 ③ 3항 조치 계획서 10-06 제출.\par
\textbf{첨부}: 1. 검수 요청 접수 확인(09-25) 2. 중간감리 1차 지적 대응 관리대장(⑭ 항목 6) 3. 귀속 이의 서면(09-29) 4. 심각도 2 미해결 목록(⑭ 항목 5) 5. 계약서 §14·§15·§18 발췌\par
\fi
\end{fieldbox}

% ------------------------------------------------------------------ 5 하도급 승인 기록
\secbar[70mm]{5. 하도급 승인 기록}
\guide{승인 요청(일자·대상 WBS·FP·업체·비율), 검토 4축(기술 / 보안·개인정보 위탁 — 재위탁 목록 반영 / 형상 — 소스 인수·SBOM 조항 승계 / 계약), 승인 조건, 승인·서명, 사후 확인(계약서 사본·재위탁 반영·투입 확인). 사후 승인을 하지 않는다. 하도급 업체의 개인정보 접근은 P4의 재위탁 구조에 반영되어야 한다(4건이 5건이 된다).}
{\footnotesize
\begin{tabularx}{\linewidth}{|L{26mm}|Y|}\hline
\hd{항목} & \hd{내용} \\ \hline
\B{\rd{10mm}}{}요청 & \Bg{[일자, 요청자, 대상 WBS·FP(전체 대비), 기간]}{2026-08-05, 한빛디지털 → P1 PM. 대상: WBS 1.2.3 배달3사 채널 연동 3본 중 연동 어댑터 개발(220 FP / 1.2.3 총 380 FP), R3 S19–S24} \\ \hline
\B{\rd{8mm}}{}업체 & \Bg{[업체명, 전문성·실적, 규모]}{㈜모빌커넥트 [추가 설정] — 배달 플랫폼 연동 전문(배달3사 API 연동 실적 7건), 12명} \\ \hline
\B{\rd{8mm}}{}하도급 비율 & \Bg{[FP × 단가 ÷ 계약금 = n\% (조항 상한, 누적)]}{220 × 55만 원 = 1.21억 / 68.2 = 1.8\% (계약 §21 상한 20\% 이내, 누적 1.8\%)} \\ \hline
\B{\rd{18mm}}{}검토 & \Bg{[기술(검토자·일자·조건) / 보안·개인정보(재위탁 해당 여부·목록 갱신) / 형상(승계 조항) / 계약(대금·연대 책임)]}{기술(박세연 08-10: 어댑터 계층 분리 구조, 소스 발주사 저장소 직접 커밋 조건) / 보안·개인정보(최영주 08-14: 주문 데이터에 고객 연락처 포함 → \textbf{개인정보 처리 재위탁 해당, 재위탁 목록 4건 → 5건 추가, P4 위탁 구조 반영·고객 고지 문구 갱신}) / 형상(PMO 08-12: 소스 인수·SBOM·라이선스 조항 하도급 계약 승계 확인) / 계약(강현도 08-17: 하도급 대금 직접 지급 요청 시 협조 조항, 수행사 연대 책임 유지)} \\ \hline
\B{\rd{14mm}}{}승인 조건 & \Bg{[① 계약서 사본 ② 재위탁 절차 완료 후 접근 ③ 조항 승계 ④ 보안 서약 ⑤ 책임 유지·교체 절차]}{① 하도급 계약서 사본 제출(08-26 접수) ② 재위탁 고지·동의 절차 P4 완료 후 개인정보 접근(09-01 완료) ③ 소스 인수·SBOM 조항 승계 ④ 하도급 인력 보안 서약·출입 등록 ⑤ 수행사 책임 유지, 하도급 인력 교체 시 항목 3-1 절차 준용} \\ \hline
\B{\rd{8mm}}{}승인 / 서명 & \Bg{[P1 PM 일자 / 구매 / 법무 / 수행사 PM]}{P1 PM 2026-08-19, 강현도(계약), 최영주(위탁 검토), 이재현} \\ \hline
\B{\rd{8mm}}{}사후 확인 & \Bg{[사본 접수 / 재위탁 반영 / 투입 확인]}{계약서 사본 08-26 완료, 재위탁 반영 09-01 완료, S19 착수 09-28 하도급 인력 4명 투입 확인} \\ \hline
\end{tabularx}}

% ------------------------------------------------------------------ 6 (가로)
\landscapeon
\secbar{6. 지체상금 · 지연 귀책 사전 정리 \B{}{— 2026-09-30, 변경협의체 10-14 상정안}}
\guide{임계경로 지연(영업일)의 원인별 분해(발주사 단독 / 수행사 단독 / 동시), 각 구간의 계약상 효과(대기 비용 / 지체상금 / 공기 연장), 여유 소유권 문안 인용과 TF 상태, 동시 지연 원칙, 지체상금 산식과 유보액, G4 정산 원칙. 지금 확정하는 것은 \textbf{귀책 분해}이고 금액은 G4에 정산한다. “복합”으로 두면 종료 시 전부 다툰다. 여유 안의 지연에는 지체상금을 말하지 않는다.}
\begin{fieldbox}\small
\B{\g{여유 소유권 문안(일정 기준선 항목 8) 인용: }\lines{2}{6mm}\g{현재 TF 상태와 동시 지연 원칙(합의 일자·조항): }\lines{2}{6mm}}{%
\textbf{여유 소유권 문안(일정 기준선 v1.0 항목 8) 재인용}: \emph{“총 여유는 발주사도 수행사도 소유하지 않으며 먼저 쓰는 쪽이 쓴다. 발주사 귀책 지연이 발생하더라도 그 시점에 해당 경로의 총 여유가 남아 있어 완료일(G4)이 밀리지 않으면 공기 연장은 없고, 총 여유가 0 아래로 떨어져 완료일이 밀리는 만큼만 연장한다(SCL 2판 원칙 8). 변경협의체는 대기 비용 판정 시 이 문안과 CPM 계산표의 TF를 근거로 한다.”}\par
\textbf{현재 상태}: R2(A06) 종료 +10영업일, A05·A06·A10은 기준선 TF 0(임계경로) → 지연 10일은 여유에 흡수되지 않고 G4를 민다. \textbf{동시 지연 원칙(SCL 원칙 10, 계약 §14·§18 해석 기준으로 변경협의체 합의 07-15)}: 발주사 원인과 수행사 원인의 지연이 같은 기간에 겹치면 수행사는 그 기간의 공기 연장(지체 제외)은 받되 대기 비용(연장 비용)은 받지 못하고, 발주사는 지체상금을 물리지 못한다.}
\end{fieldbox}
{\footnotesize\setlength{\tabcolsep}{2.5pt}
\begin{tabularx}{\linewidth}{|L{46mm}|C{18mm}|Y|L{34mm}|L{42mm}|L{56mm}|}\hline
\hd{구간} & \hd{지연 (영업일)} & \hd{원인} & \hd{귀책} & \hd{계약상 효과} & \hd{판정 · 유보} \\ \hline
\ifwbblank
\rd{8mm}\g{[스프린트·활동, 기간]} & & \g{[결정 SLA 초과분 등 — 사실로]} & \g{[발주 단독 / 수행 단독 / 동시]} & \g{[대기 비용 · 지체 산입 · 공기 연장]} & \g{[판정 완료(일자·서면) / 유보 — G4]} \\ \hline
\rd{8mm} & & & & & \\ \hline
\rd{8mm} & & & & \g{[겹침 → 동시]} & \\ \hline
\rd{8mm} & & & & & \\ \hline
\rd{8mm}\g{[회복 구간]} & \g{[−n]} & \g{[속도·초과 근무 — 부담 주체]} & — & — & \g{[상계 주장 가능 여부]} \\ \hline
\rd{8mm}\g{[병행 가능 지연]} & \g{[n (임계 영향 0)]} & & & \g{[해당 부분 한정]} & \\ \hline
\rd{8mm}\textbf{순 지연 (임계경로)} & \g{[n]} & & \g{[발주 n + 동시 n + 수행 n − 회복 n]} & & \g{[G4 영향]} \\ \hline
\else
S10 (05-11~05-29) 3주 연장 & 7 & POS 화면 정의서 승인 대기 12영업일 중 결정 SLA(5일) 초과분 & \textbf{발주사 단독} & 대기 비용 대상, 지체 제외 & \textbf{판정 완료 07-15} — 0.42억 지급, 합의서 5호 \\ \hline
S10 & 3 & 승인 후 재작업 중 수행사 결함(POS 할인 로직) 재시험 & \textbf{수행사 단독} & §18 지체 일수 산입 & 유보 — G4 정산 \\ \hline
S18 (09-07~09-25) 3주 연장 & 2 & 규제 검증 중 최영주 추가 해석 의견(파기 주기 로그, 09-01) 반영 & 발주사 & 겹침 → 동시 & 동시 지연 3일로 합산 \\ \hline
S18 & 3 & R2 결함 마무리(심각도 2 이상 9건) & 수행사 & 겹침 → 동시 & \textbf{동시 지연: 공기 연장 인정, 대기 비용·지체상금 불인정} (변경협의체 10-14 확인 예정) \\ \hline
S12–S17 회복 & −3 & 속도 640(초과 근무) — 수행사 부담 & — & — & 회복분은 수행사 단독 귀책 3일에서 우선 상계 주장 가능 — 10-14 협의 \\ \hline
AF-1-03·04 귀속 정리 (09-28~10-08) & 8 (병행 가능, 임계경로 영향 0) & 발주사 결정 미완 & \textbf{발주사 단독} (영향 범위 R2 FP 12\%) & 대기 비용 대상(해당 부분 한정), 지체 제외 & 대기 기록 10-09 제출 후 10-14 판정 \\ \hline
\textbf{순 지연 (임계경로)} & \textbf{10} & & 발주 7 + 동시 3 + 수행 3 − 회복 3 & & G4 영향 +10(회복 계획 EX-01로 0 목표) \\ \hline
\fi
\end{tabularx}}
\B{\small\g{지체상금 산정 예시(유보액 — 지연 부분 계약금의 정의 포함) / G4 정산 원칙 ①~④ / 공기 연장(EOT) 청구 현황: }\lines{3}{6mm}}{\small \textbf{지체상금 산정 예시(유보)}: 수행사 단독 3일 × 1/1000 × 지연 부분 계약금(R2 FP 3,720 × 55만 원 = 20.46억) = \textbf{0.0614억(614만 원)} — 회복 계획 성공 시 G4 완료일이 기준선과 같아지므로 §18의 “지연”이 성립하지 않아 소멸 가능. \textbf{G4 정산 원칙}: ① 최종 완료일 기준 순 지연 일수 확정 ② 귀책 분해표(본표 누적) 적용 ③ 동시 구간 상계 ④ 대기 비용 기지급분 정산. \textbf{공기 연장(EOT) 청구}: 수행사 09-30 현재 없음(회복 계획 동의). 검산: 3,720 × 0.0055 = 20.46; 20.46 × 3/1000 = 0.06138.\par}

% ------------------------------------------------------------------ 7 계약 변경 로그 (가로)
\clearpage
\secbar{7. 계약 변경 로그 \B{}{— 2026-01-05 ~ 09-30, 예정 포함 10건}}
\guide{모든 계약의 변경 한 줄 — 번호·일자·계약·내용·금액 변동·사유(CR·R·IS 연결)·승인자·서면 번호·상태. P1 예산 밖 계약도 소관 표시로 포함. \textbf{금액 0인 변경(조건 변경)도 적는다} — 지급 조건 50/50이 EVM 착시의 원인이었다. 누적 행은 P1 예산 내 금액 변동.}
{\scriptsize\setlength{\tabcolsep}{2pt}
\begin{longtable}{|C{8mm}|L{18mm}|L{34mm}|L{68mm}|L{30mm}|L{22mm}|L{32mm}|L{22mm}|L{20mm}|}\hline
\hd{\#} & \hd{일자} & \hd{계약} & \hd{내용} & \hd{금액 Δ (억)} & \hd{사유 (CR/R/IS)} & \hd{승인자} & \hd{서면} & \hd{상태} \\ \hline \endhead
\ifwbblank
\rd{5mm}1 & & & & \g{[0이면 “지급 시점 이동 n”]} & & & & \\ \hline
\rd{5mm}2 & & & & & & & & \\ \hline
\rd{5mm}3 & & & & \g{[유보/이관 표시]} & & & & \\ \hline
\rd{5mm}4 & & & & & & & & \\ \hline
\rd{5mm}5 & & \g{[P1 예산 외 — 소관]} & & \g{[0 (P1 예산 외)]} & & & & \\ \hline
\rd{5mm}6 & & & & & & & & \\ \hline
\rd{5mm}7 & & & & & & & & \\ \hline
\rd{5mm}8 & & & & & & & & \\ \hline
\rd{5mm}9 & & & & & & & & \\ \hline
\rd{5mm}10 & \g{[예정]} & & & & & & & \g{[예정]} \\ \hline
\rd{5mm}\textbf{누적 (P1 예산 내)} & & & & \g{[계약별 순변동]} & & & & \\ \hline
\else
1 & 03-13 & 상용SW 3사 (MDM 2차·APM, SAP 애드온) & 지급 조건 “인도 시 100\%” → “인도 50\% / 통합시험 통과 50\%”, 인도 5월·7월로 조정 & 0 (지급 시점 이동 4.70) & IS-01 / R-04 & 정미란(03-06), 강현도 & SW-01·02·03 & 완료 \\ \hline
2 & 03-25 & SI & CR-03 엑셀 양식 5종 +22 FP & +0.121 & CR-03 & P1 PM 전결 & SI 합의서 3호 & 완료 \\ \hline
3 & 03-27 → 04-24 & 클라우드 & 3년 선약정 14.6 → 2년 선약정 6.0 + 종량(기준선 11.6), 차액 3.0 유보 & −3.0 (유보) & Day2 조정표 & 박세연·강현도, 정미란 & CL-01 & 완료 (04-24 서명) \\ \hline
4 & 04-10 & 상용SW 유지보수 & 2·3년차 유지보수 4.0 P5 운영비 계약으로 분리 & −4.0 (P5 이관) & Day2 조정표 & 박준영·재경팀장, 정미란 & SW-04 & 완료 \\ \hline
5 & 04-16 & SI & CR-04 옴니 고도화 3종 R5 이연 확인(FP 불변, 인도 시점) & 0 & CR-04 & CCB & SI 합의서 4호 & 완료 \\ \hline
6 & 04-22 & \emph{앱 벤더 유지보수 (영업본부)} & 소스 인수·API 연동 시험 협조 조항 추가 & 0 (P1 예산 외) & R-08 & 오정근·박세연 & 앱-01 & 완료 — R-18 주시 \\ \hline
7 & 07-08 & \emph{3PL (물류본부)} & 준실시간 연동 불가 확인, 야간 배치 2회 폴백, 2027-06 재협상 조항 & 0 (P1 예산 외) & IS-05 / R-03 & 한동석 & 3PL 부속 & 완료 \\ \hline
8 & 07-15 & SI & CR-06 +58 FP + R-05 대기 비용 0.42 & +0.739 & CR-06, IS-02 & 변경협의체 → P1 PM·CCB 사후 & SI 합의서 5호 & 완료 (대기 비용 07-31 지급) \\ \hline
9 & 08-19 & SI (하도급) & 배달3사 연동 어댑터 220 FP 하도급 승인(모빌커넥트) & 0 & 항목 5 & P1 PM, 강현도, 최영주 & 하도급 승인서 1호 & 완료 \\ \hline
9a & 09-01 & 감리 & 중간감리를 2회(1차 09월 설계·R2, 2차 G2 11월)로 분할 실시 — 금액 1.4 불변(0.7/0.7) & 0 & 감리 계획 & P1 PM, 정하늘 & AU-01 & 완료 \\ \hline
10 & 10-23 (예정) & SI + 데이터 전환 용역 & SI 합의서 6호: CR-07·08·09·11 (+40 +90 +35 −60 = +105 FP) / DT-01: 대상 5년 한정 11.4 → 9.4 & SI +0.5775 / 용역 −2.0 & CCB 10-15 & 강현도, 정미란 & SI 6호, DT-01 & 예정 \\ \hline
\textbf{누적 (P1 예산 내)} & & & & SI \textbf{+1.4375}(FP +185 · 대기 0.42) / 용역 −2.0 / 클라우드 −3.0(유보) / 유지보수 −4.0(이관) / 상용SW 시점 이동 4.70 & & & & \\ \hline
\fi
\end{longtable}}
\B{}{\small 검산: SI 합의서 3호 0.121 + 5호 0.739 + 6호 0.5775 = 1.4375 (= FP 185 × 0.0055 + 0.42). 6호 FP: +40 +90 +35 −60 = +105, 105 × 0.0055 = 0.5775. 변경 로그의 CR 번호는 ⑪ 항목 10과 같아야 한다.\par}

% ------------------------------------------------------------------ 8 조달 상태 (가로)
\secbar[70mm]{8. 조달 상태와 다음 조달}
\guide{미완 조달(단말·DR 장비·부하시험 3자·계약 내 잔여 회차)의 상태·견적(기준선 대비)·발주 → 입고·완료(리드타임 역산)·결재선(5억 초과 스폰서 병행, 5,000만 원 이상 CFO)·리스크(등록부 연결). 조달 리드타임을 일정에 넣지 않으면 12월 봉우리가 1월로 밀린다.}
{\footnotesize\setlength{\tabcolsep}{2.5pt}
\begin{tabularx}{\linewidth}{|L{34mm}|L{52mm}|L{26mm}|Y|L{38mm}|L{50mm}|}\hline
\hd{조달} & \hd{상태 (데이터 일자)} & \hd{견적 · 금액 (억)} & \hd{발주 → 입고 · 완료 (리드타임)} & \hd{결재선} & \hd{리스크} \\ \hline
\ifwbblank
\rd{12mm}\g{[물품]} & \g{[견적 접수 일자, 기준선 대비]} & & \g{[입고 필요일 ← 리드타임 = 발주 기한]} & \g{[금액 구간별 결재]} & \g{[지연 시 영향 월 · 등록부 R-]} \\ \hline
\rd{12mm}\g{[장비]} & & & & & \\ \hline
\rd{12mm}\g{[3자 용역]} & & & & & \\ \hline
\rd{10mm}\g{[계약 내 잔여 회차]} & & & & — & \\ \hline
\rd{10mm} & & & & & \\ \hline
\else
POS 단말 477대 + 주변기기 & 단일 공급사 견적 접수 09-18: 8.9(기준선 8.6, 반도체 단가 +0.3 [추가 설정]) & 8.9 & 발주 10-30 → 입고 12-14(리드타임 6주) → 설치 2027-01~02(설 잠금 01-24~02-06 회피) & 5억 초과 → 정미란 병행 결재 10-23, 내부통제 지침 & 입고 지연 시 2027-01 봉우리(단말 4.86) 이월 → R-17과 결합 \\ \hline
DR·백업 장비 & 사양 확정 09-25(박세연) & 2.1 + 2.1 & 발주 11-13 → 12월·1월 설치 & 5,000만 원 이상 CFO 결재 & 클라우드 DR 구성으로 대체 검토(IS-10 사이징과 연계) \\ \hline
부하시험 3자 & RFP 초안 & 1.5 + 1.6 + 시나리오 추가 0.2 & 계약 11-27 → 2027-01 실시 & 5,000만 원 이상 & 피크 5,100 케이스/h 시나리오 — 유보 3.0 복귀 판단 근거 \\ \hline
보안진단 2회차 & 계약 내 & 1.4 & 2027-01 & — & High 0 조건(G3) \\ \hline
감리 G2 2차·G3·G4 & 계약 내 & 0.7 + 1.4 + 0.6 & 11월·2월·5월 & — & 지적 귀속 판정 원칙 합의(AF-1-03·04 회신 결과) \\ \hline
\fi
\end{tabularx}}
\B{\small\g{서명: P1 PM \hrulefill\ 구매팀장 \hrulefill\ 재경팀장 \hrulefill\ 스폰서(5억 초과 건) \hrulefill}\par}{\small 서명(10-05): P1 PM / 강현도 / 재경팀장 / 정미란(단말 발주 병행 결재는 10-23). 변경협의체 10-14 상정: 항목 4 대기 판정, 항목 6 동시 지연 확인, 항목 2 합의서 6호 초안.\par}
\landscapeoff
\end{document}
